% page 273 of the facsimile (image p-272 of the scan)
% block 162.6 x 242.0 mm ; left margin 39.4 mm
\pgc{17.780mm}{0.000mm}{
\l{}
\l{}
\l{\cel{51}265}
\l{}
\l{\sou{karcinomo}. (\sou{patol}.)Un ek la formi di la morbo \textquotedbl{}kancero\textquotedbl{}.DEFIRS.}
\l{}
\l{\sou{kardar}.\cel{2}(\sou{trans.}) Pektar, desintrikar (precipue lano e krino)}
\l{\cel{3}per planketo provizita ye klovi ek latuno-filo. - DEFIRS.}
\l{}
\l{\sou{kardamino}. (\sou{bot.)}\cel{3}Planto krucifera. - L.\cel{1}\sou{cardamina}. - DEFIS.}
\l{}
\l{\sou{kardamomo}. (\sou{bot.}) Planto di India, genero \textquotedbl{}amomo\textquotedbl{}. - L.\cel{1}\sou{elitte}-}
\l{\cel{3}\sou{ria cardamomum}. - DEFIRS.}
\l{}
\l{\sou{kardano}. (\sou{tekn.)}\cel{2}\textquotedbl{}Artiko\textquotedbl{} di parti de mekanismi (stangi, e c.)}
\l{\cel{3}qua posibligas movi irga-direcione. - DEFIRS.}
\l{}
\l{\sou{kardelo}. (\sou{zool.}) Mikra ucelo kantera, de la serio \textquotedbl{}paseri\textquotedbl{}, di}
\l{\cel{3}qua la maskulo havas kapo reda e la ali markizita flave e}
\l{\cel{3}brune. - L.\cel{1}\sou{garduelis elegans}. - FIL.}
\l{}
\l{kardinala. I. (\sou{matem.)}\cel{2}Dicesas pri la nombri qui indikas quanta}
\l{\cel{3}unaji sama-natura kontenesas en quanto determinita, e qui}
\l{\cel{3}uzesas por formacar la cetera termini numerala. - II. Dic-}
\l{\cel{3}esas pri la quar horizonto-punti : nordo, sudo, esto, westo,}
\l{\cel{3}segun qui onu determinas la situeso di loki cetera. - III.}
\l{\cel{3}Dicesas pri la quar vertui precipua (prudenteso, yusteso,}
\l{\cel{3}moderemeso, e psiko-forteso). DEFIS.}
\l{}
\l{kardinalo.\cel{2}Singla ek la sepa-dek prelati (episkopi, sacerdoti,}
\l{\cel{3}diakoni) ye qui konsistas la Sakra Kolegio di la Eklezio}
\l{\cel{3}katolika.qui votas en konklavo por elektar nova papo, o por}
\l{\cel{3}elektesar, e qui havas kostumo reda. - DEFIRS.}
\l{}
\l{kardioido. (\sou{geom.)}\cel{2}Konkoido di la cirklo.}
\l{}
\l{kardono. (\sou{bot.}) Planto ek la familio \textquotedbl{}kompozaji\textquotedbl{}, di qua la}
\l{\cel{3}folii e la kapituli esas dornoza. - L.\cel{1}\sou{carduus}. - EFIS.}
\l{}
\l{karduno. (\sou{bot}.)\cel{2}Speco di artichoko, di qua onu manjas la}
\l{\cel{3}kosto mediana di la folii, pos blankigir olu per kovrar lu}
\l{\cel{3}ye palio. - L.\cel{1}\sou{cynara cardunculus}. - DEFIS.}
\l{}
\l{karear. (\sou{trans.}) Ne judikar su kom privacata de ulo, pro ke}
\l{\cel{3}onu ne tre bezonas nek tre deziras olu. - EFL.}
\l{}
\l{karelo. I. (\sou{ludo-karto)}. Quadrato reda, uzata kom la distingivo}
\l{\cel{3}di grupo de karti, en qua la aso reprezentesas per quadrato}
\l{\cel{3}meze di la karto. - II. Objekto di qua la surfaco esas qua-}
\l{\cel{3}drata o rektangula. - FL.}
\l{}
\l{kareno. La tota surfaco di navo,qua esas imersita. - FIS.}
\l{}
\l{\cel{1}\sou{karesmo}. En la religio katolika, periodo ek 46 dii, de la cin-}
\l{\cel{4}dro-merkurdio til la Pasko-sundio,dum qua la fideli devas}
\l{\cel{4}abstinencar e fastar (ecepte dum la sundii). - FIS.}
}
